\section{Maximum Path Sum}
\label{sec:trees-max-path-sum}

\problemstatement{A path is any sequence of nodes connected by edges, where
each adjacent pair shares an edge and no node repeats; it need not pass
through the root. The path sum is the total of its node values. Given the
root (values may be negative), return the maximum path sum.}

\begin{patternbox}
This is the diameter problem with \emph{values instead of edge counts},
plus one subtlety: because values can be negative, a subtree that would
\emph{lower} the total should be dropped entirely. The keywords
\emph{``maximum,'' ``path,'' ``may be negative''} say: bottom-up, track a
global best through each node, but clamp negative child contributions to
zero.
\end{patternbox}

\subsection*{Understanding the problem carefully}

At any node, the best path with that node as its peak is: the node's value,
plus the best downward path into the left subtree (if positive), plus the
best downward path into the right subtree (if positive). That is the
candidate we compare against a global best. But the value we \emph{return}
to the parent can only continue \emph{one} way -- a path cannot fork at the
parent -- so we return the node's value plus the better of its two
downward gains.

\begin{ideabox}[Return a Straight-Down Gain, Score a Bent Path]
Let \textsc{Gain}(node) return the maximum sum of a path that starts at
\textit{node} and goes strictly downward. Compute
$l=\max(0,\textsc{Gain}(\text{left}))$ and
$r=\max(0,\textsc{Gain}(\text{right}))$ (the $\max(0,\cdot)$ drops
harmful subtrees). Update the global best with $node.val+l+r$ (the path
bending through this node). Return $node.val+\max(l,r)$ to the parent (only
one branch may continue).
\end{ideabox}

\subsection*{Algorithm}

\begin{enumerate}
  \item Global \textit{best} $=-\infty$.
  \item \textsc{Gain}(\textit{node}):
  \begin{enumerate}
    \item If null, return $0$.
    \item $l=\max(0,\textsc{Gain}(\text{left}))$;
          $r=\max(0,\textsc{Gain}(\text{right}))$.
    \item \textit{best} $=\max(\textit{best},\ node.val+l+r)$.
    \item Return $node.val+\max(l,r)$.
  \end{enumerate}
  \item Answer is \textit{best}.
\end{enumerate}

\subsection*{Dry run}

\dryrun{Tree: root $-10$; left $9$; right $20$ with children $15,7$.}

\begin{itemize}
  \item Gain($9$)$=9$, Gain($15$)$=15$, Gain($7$)$=7$.
  \item At $20$: $l=15,r=7$; best$=\max(-\infty,20+15+7)=42$; return
        $20+15=35$.
  \item At $-10$: $l=\max(0,9)=9$, $r=\max(0,35)=35$;
        best$=\max(42,-10+9+35)=42$; return $-10+35=25$.
  \item Maximum path sum $=42$ (path $15\!-\!20\!-\!7$).
\end{itemize}

\subsection*{C++ implementation}

\begin{lstlisting}
#include <bits/stdc++.h>
using namespace std;

struct TreeNode {
    int val;
    TreeNode* left;
    TreeNode* right;
    TreeNode(int v) : val(v), left(nullptr), right(nullptr) {}
};

int best;

int gain(TreeNode* node) {
    if (node == nullptr) return 0;
    int l = max(0, gain(node->left));   // drop harmful subtrees
    int r = max(0, gain(node->right));
    best = max(best, node->val + l + r); // path bending through node
    return node->val + max(l, r);        // parent can continue one way
}

int maxPathSum(TreeNode* root) {
    best = INT_MIN;
    gain(root);
    return best;
}

int main() {
    TreeNode* root = new TreeNode(-10);
    root->left = new TreeNode(9);
    root->right = new TreeNode(20);
    root->right->left = new TreeNode(15);
    root->right->right = new TreeNode(7);

    cout << "Max path sum: " << maxPathSum(root) << "\n";
    return 0;
}
\end{lstlisting}

\begin{complexitybox}
\textbf{Time Complexity.} $O(n)$ -- one postorder pass. \textbf{Space
Complexity.} $O(h)$ for the recursion stack.
\end{complexitybox}

\begin{takeawaybox}
The $\max(0,\cdot)$ clamp -- ``take a subtree's contribution only if it
helps'' -- is the one new idea beyond the diameter template. Whenever a
downward contribution can be negative, clamping it to zero is how greedy
``skip the bad branch'' decisions are expressed inside a bottom-up
recursion.
\end{takeawaybox}
